\item \textbf{Distribución de estrellas en una galaxia (modelo simple).}

\begin{tcolorbox}[colback=blue!2!white,colframe=blue!55!black,title={Enunciado}]

Considere una galaxia idealizada donde las estrellas están distribuidas en el plano \((x,y)\) siguiendo una distribución aproximadamente gaussiana centrada en el origen.

Se desea generar un conjunto de posiciones de estrellas y calcular magnitudes asociadas a su posición.

\begin{itemize}
\item[(a)] Genere \(N = 1000\) valores para el eje \(x\) y \(N\) valores para el eje \(y\) usando la función \texttt{numpy.random.normal}, con media \(0\) y desviación estándar \(\sigma = 100 \).

\item[(b)] Considere que cada par \((x_i, y_i)\) representa la posición de una estrella en la galaxia.

\item[(c)] Calcule la distancia de cada estrella al centro usando:
\[
r_i = \sqrt{x_i^2 + y_i^2}.
\]

\item[(d)] Determine:
\begin{itemize}
\item la distancia máxima al origen,
\item la distancia mínima,
\item y el promedio de las distancias.
\end{itemize}

\end{itemize}

\end{tcolorbox}

\begin{solution}

\subsubsection*{(a) Generación de posiciones}

La función \texttt{numpy.random.normal} permite generar números aleatorios con distribución normal. En este caso, generamos dos arreglos independientes:

\begin{pythonjupyter}
import numpy as np

# Número de estrellas
N = 1000

# Parámetros de la distribución
sigma = 1.0

# Generación de posiciones
x = np.random.normal(0, sigma, N)
y = np.random.normal(0, sigma, N)
\end{pythonjupyter}

Cada posición \((x_i, y_i)\) representa una estrella en el plano.

\subsubsection*{(b) Interpretación de los datos}

Los arreglos \texttt{x} e \texttt{y} contienen las coordenadas de las estrellas. Ambas componentes son independientes y están centradas en el origen.

\subsubsection*{(c) Cálculo de la distancia al origen}

Se utiliza la definición de distancia euclidiana:

\begin{pythonjupyter}
# Distancia al origen
r = np.sqrt(x**2 + y**2)
\end{pythonjupyter}

Esta operación se realiza de forma vectorizada sobre todos los elementos del arreglo.

\subsubsection*{(d) Magnitudes globales}

Se calculan valores característicos del arreglo de distancias:

\begin{pythonjupyter}
# Estadísticas básicas
r_max = np.max(r)
r_min = np.min(r)
r_mean = np.mean(r)

print("Distancia máxima:", r_max)
print("Distancia mínima:", r_min)
print("Distancia promedio:", r_mean)
\end{pythonjupyter}

Estos valores permiten describir de forma global la distribución espacial de las estrellas.

\end{solution}
